% benötigt im Präambel: \usepackage{booktabs,tabularx,amsmath}
\begin{table}[htbp]
  \centering
  \caption[Kennzahlenblatt]{Kennzahlenblatt (vgl. zur Ausgestaltung von Kennzahlenblättern z.\,B. Schulte 2002, S.~170, und Eisele und Doyé 2010, S.~384)}
  \label{tab:4-8}
  \small
  \begin{tabularx}{\textwidth}{@{}>{\raggedright\arraybackslash}p{3.3cm}>{\raggedright\arraybackslash}X@{}}
    \toprule
    \multicolumn{2}{@{}l}{\textbf{Kennzahl: Anlagendeckung (\%)}} \\
    \midrule
    Definition & Die Anlagendeckung gibt darüber Auskunft, inwieweit das Anlagevermögen durch langfristiges Kapital (Eigenkapital + langfristiges Fremdkapital) gedeckt ist. \\
    \addlinespace
    Formel & $\displaystyle \text{Anlagendeckung (\%)} = \frac{\text{Eigenkapital} + \text{langfristiges Fremdkapital}}{\text{Anlagevermögen}} \cdot 100$ \\
    \addlinespace
    Einheit & Prozentwert \\
    Zielgröße & 110--150\,\% \\
    Wertearten & Ist, Plan, Forecast \\
    Basisdaten und Datenquelle & Bilanzkennzahlen: Eigenkapital, langfristiges Fremdkapital und Anlagevermögen (aus der Finanzbuchhaltung) \\
    Objektzuordnungen & Gesellschaft \\
    Anwendungsbereiche & Spitzenkennzahlen Bilanz-Controlling \\
    Empfänger/ Rollenzuordnung & Geschäftsführung, Führungsebene \\
    Erhebungszyklus & Quartal, Jahr \\
    \addlinespace
    Bedeutung & Im Rahmen der Beobachtung der langfristigen Finanzstruktur spielt die Kennzahl Anlagendeckung eine wichtige Rolle. Das dieser Kennzahl zugrunde liegende Prinzip ist die fristenkongruente Finanzierung des langfristig gebundenen Vermögens, das über die Anlagendeckung kontrolliert werden soll und die Funktion hat, langfristig strukturelle Ungleichgewichte aufzudecken.

    Langfristiges Vermögen soll auch langfristig finanziert sein (goldene Bilanzregel). Deshalb sollte der Deckungsgrad~II deutlich über 100\,\% liegen (Ziel z.\,B. 150\,\%). Je weiter der Deckungsgrad~II über 100\,\% liegt, umso mehr ist neben dem Anlagevermögen auch das Umlaufvermögen durch langfristiges Kapital finanziert und damit eine höhere finanzielle Stabilität des Unternehmens gegeben. Ist das Anlagevermögen z.\,B. zum Teil kurzfristig finanziert, könnte das Unternehmen bei Fälligkeit kurzfristiger Verbindlichkeiten in Zahlungsschwierigkeiten geraten, da das Umlaufvermögen nicht ausreicht und das Anlagevermögen nicht so schnell liquidierbar ist. \\
    \bottomrule
  \end{tabularx}
\end{table}
